\chapter{Representações de grupos finitos}\label{ch:b3:representations}

Para compreender um grupo abstrato, faça-o agir sobre um espaço
vetorial e aplique álgebra linear --- autovalores, traços, produtos
internos --- à ação. Esse programa, a \emph{teoria de \omterm{def:b3:representations:rep}{representações}},
é espantosamente eficaz para grupos finitos sobre $\C$: toda
\omterm{def:b3:representations:rep}{representação} se decompõe em irredutíveis (Maschke), as irredutíveis
ficam determinadas por seus \emph{\omterm{def:b3:representations:character}{caracteres}} (traços), e os \omterm{def:b3:representations:character}{caracteres}
satisfazem relações de ortogonalidade que tornam mecânicos os cálculos.
O capítulo constrói esse cálculo e suas primeiras obras-primas --- as
\omterm{def:b3:representations:table}{tabelas de caracteres} dos grupos pequenos ---, e o problema de fim de
semana colhe um teorema muito além do alcance da teoria de grupos nua:
o teorema $p^aq^b$ de Burnside. Ao longo de todo o capítulo, $G$ é
um grupo finito e todos os espaços vetoriais têm dimensão finita sobre
$\C$.

\section{Representações, Maschke, Schur}

\begin{definition}\label{def:b3:representations:rep}
Uma \emph{representação}\index{representação} de $G$ é um morfismo
$\rho \colon G \to GL(V)$ para um $\C$-espaço vetorial $V$; $\dim V$ é seu
\emph{grau}. Um subespaço $W \subseteq V$ é \emph{invariante} se $\rho(g)W \subseteq W$ para todo
$g$; a restrição faz de $W$ uma \emph{subrepresentação}.
$\rho$ é \emph{irredutível}\index{representação irredutível} se
$V \neq
0$ e seus únicos subespaços invariantes são $0$ e $V$. Um
\emph{morfismo} entre $(\rho, V)$ e $(\sigma, W)$ é uma aplicação linear $f\colon V \to W$ com $f\rho(g) = \sigma(g)f$
para todo $g$ (equivariância); os $f$ bijetores são
\emph{isomorfismos}.
\end{definition}

\begin{example}\label{ex:b3:representations:examples}
(a) Grau $1$: morfismos $G \to \C^\times$.
(b) A \emph{representação regular}\index{representação regular}: $V = \C^G$
com base $(e_h)_{h \in G}$, $\rho(g)e_h = e_{gh}$; de grau $\abs G$.
(c) Uma ação por permutações de $G$ sobre um conjunto finito $X$
dá a \emph{\omterm{def:b3:representations:rep}{representação} de permutação} em $\C^X$: $\rho(g)e_x =
e_{g\cdot x}$.
(d) $S_n$ age em $\C^n$ permutando as coordenadas; o hiperplano
$\{\sum x_i = 0\}$ é invariante: é a \emph{\omterm{def:b3:representations:rep}{representação} padrão}, de grau
$n - 1$.
\end{example}

\begin{theorem}[Maschke]\label{thm:b3:representations:maschke}
Todo subespaço invariante $W$ de uma \omterm{def:b3:representations:rep}{representação} $(V, \rho)$ admite um
complemento invariante. Consequentemente, toda \omterm{def:b3:representations:rep}{representação} é uma soma
direta de irredutíveis (\emph{semissimplicidade}).\index{teorema de
Maschke}
\end{theorem}

\begin{proof}
Seja $p \colon V \to V$ \emph{uma} projeção qualquer de imagem $W$ (escolha um
complemento qualquer). Faça a média sobre o grupo:
\[
\tilde p = \frac1{\abs G}\sum_{g \in G} \rho(g)\,p\,\rho(g)^{-1}.
\]
Cada parcela leva $V$ em $W$ ($W$ é invariante) e fixa $W$
ponto a ponto: para $w \in W$, $\rho(g)^{-1}w \in W$, $p$ o fixa, e $\rho(g)$ o traz de volta
--- de modo que $\tilde p$ é novamente uma projeção sobre $W$. Ela é
equivariante: para $h \in G$, $\rho(h)\tilde p\rho(h)^{-1}$ apenas reindexa a mesma soma. Logo $\ker
\tilde p$
é um complemento invariante de $W$. Iterando nas parcelas (dimensão
finita), decompõe-se $V$ em irredutíveis.
\end{proof}

\begin{theorem}[Lema de Schur]\label{thm:b3:representations:schur}
Seja $f \colon V \to W$ um morfismo de \omterm{def:b3:representations:rep}{representações} \emph{irredutíveis}. Então
$f = 0$ ou $f$ é um isomorfismo; e, se $(V,\rho) = (W,\sigma)$, então $f = \lambda\,\mathrm{id}$ para
algum $\lambda \in \C$. Logo $\dim\operatorname{Hom}_G(V, W)$ vale $1$ se $V \cong W$ e $0$ caso
contrário.\index{lema de Schur}
\end{theorem}

\begin{proof}
$\ker f$ e $\operatorname{im} f$ são invariantes (equivariância), de modo que cada um
é $0$ ou o espaço todo: ou $f = 0$, ou $f$ é injetora com
imagem total. Se $V = W$: $f$ tem um autovalor $\lambda$ ($\C$ é
\omterm{def:b3:galois:closure}{algebricamente fechado}); $f - \lambda\,\mathrm{id}$ é um morfismo $V \to V$ não injetor, logo é
$0$. Quanto à contagem de dimensão quando $V \cong W$: fixando um
isomorfismo $u$, todo morfismo $f$ dá o endomorfismo $u^{-1}f = \lambda\,\mathrm{id}$: $f =
\lambda u$.
\end{proof}

\section{Caracteres e ortogonalidade}

\begin{definition}\label{def:b3:representations:character}
O \emph{caráter}\index{caráter} de $(V, \rho)$ é $\chi_\rho(g) = \operatorname{tr}\rho(g)$. Ele satisfaz $\chi_\rho(e) = \dim V$,
$\chi_\rho(hgh^{-1}) = \chi_\rho(g)$ (os traços são invariantes por conjugação): os caracteres são
\emph{funções de classe}\index{função de classe} --- elementos do
espaço $\mathcal{CF}(G)$ das funções constantes nas \omterm{ex:b3:groups:actions}{classes de conjugação}, munido do
produto interno hermitiano
\[
\langle \varphi, \psi\rangle = \frac1{\abs G}\sum_{g \in G}
\overline{\varphi(g)}\,\psi(g).
\]
\end{definition}

\begin{proposition}\label{prop:b3:representations:charbasics}
$\rho(g)$ é diagonalizável com autovalores raízes da unidade; $\chi_\rho(g^{-1}) = \overline{\chi_\rho(g)}$, e $\abs{\chi_\rho(g)} \leq \chi_\rho(e)$
com igualdade se, e somente se, $\rho(g)$ é escalar. Os \omterm{def:b3:representations:character}{caracteres} se somam
nas somas diretas: $\chi_{V
\oplus W} = \chi_V + \chi_W$.
\end{proposition}

\begin{proof}
$\rho(g)^N = \mathrm{id}$ para $N = \abs G$ (Lagrange): $\rho(g)$ aniquila $X^N - 1$ e se decompõe com
raízes simples, de modo que o operador é diagonalizável com autovalores
$\lambda_i \in \mu_N$. Então $\chi(g^{-1}) = \sum \lambda_i^{-1} = \sum\bar\lambda_i =
\overline{\chi(g)}$, e $\abs{\chi(g)} = \abs{\sum\lambda_i}
\leq \dim V$, com igualdade na desigualdade triangular
se, e somente se, todos os $\lambda_i$ são iguais, isto é, $\rho(g) = \lambda\,\mathrm{id}$. A
aditividade nas somas diretas: traços por blocos.
\end{proof}

\begin{lemma}\label{lem:b3:representations:homtrace}
Sejam $(V, \rho)$, $(W, \sigma)$ \omterm{def:b3:representations:rep}{representações}. O operador de média em $\operatorname{Hom}(W, V)$,
\[
c(A) = \frac1{\abs G}\sum_{g}\rho(g)\,A\,\sigma(g)^{-1},
\]
é uma projeção sobre $\operatorname{Hom}_G(W, V)$ e, para a aplicação $\Phi_g \colon A \mapsto \rho(g)A\sigma(g)^{-1}$, tem-se $\operatorname{tr}\Phi_g =
\chi_\rho(g)\,\overline{\chi_\sigma(g)}$.
\end{lemma}

\begin{proof}
$c(A)$ é equivariante (reindexe a soma como em Maschke), e $c$
fixa as aplicações equivariantes (cada parcela vale $A$): $c$ é
uma projeção de imagem $\operatorname{Hom}_G(W, V)$. Traço: em bases, $\Phi_g(A) = BAC$ com $B = \rho(g)$, $C = \sigma(g)^{-1}$;
na base $(E_{kl})$ das matrizes, $BE_{kl}C = \sum_{m,n}
b_{mk}c_{ln}E_{mn}$, de modo que o coeficiente de
$E_{kl}$ em $\Phi_g(E_{kl})$ é $b_{kk}c_{ll}$: $\operatorname{tr}\Phi_g =
\sum_{k,l}b_{kk}c_{ll} = \operatorname{tr}(B)
\operatorname{tr}(C) = \chi_\rho(g)\chi_\sigma(g^{-1})$ e $\chi_\sigma(g^{-1}) = \overline{\chi_\sigma(g)}$.
\end{proof}

\begin{theorem}[Primeiras relações de ortogonalidade]
\label{thm:b3:representations:orthogonality}
Sejam $\rho, \sigma$ \emph{irredutíveis}. Então\index{ortogonalidade dos
caracteres}
\[
\langle\chi_\sigma, \chi_\rho\rangle =
\begin{cases}
1 & \text{se } \rho \cong \sigma,\\
0 & \text{caso contrário:}
\end{cases}
\]
os \omterm{def:b3:representations:character}{caracteres} irredutíveis formam uma família ortonormal em $\mathcal{CF}(G)$.
\end{theorem}

\begin{proof}
O traço de uma projeção é a dimensão de sua imagem:
\[
\dim\operatorname{Hom}_G(W, V) = \operatorname{tr} c
= \frac1{\abs G}\sum_g \operatorname{tr}\Phi_g
= \frac1{\abs G}\sum_g \chi_\rho(g)\overline{\chi_\sigma(g)}
= \langle \chi_\sigma, \chi_\rho\rangle,
\]
e o lema de Schur avalia o membro esquerdo em $\delta_{\rho \cong
\sigma}$.
\end{proof}

\begin{corollary}\label{cor:b3:representations:multiplicity}
Decomponha $V \cong \bigoplus_i V_i^{\oplus m_i}$ em irredutíveis distintas ($V_i \not\cong V_j$). Então $m_i = \langle
\chi_{V_i}, \chi_V\rangle$: as
multiplicidades --- e portanto a \omterm{def:b3:representations:rep}{representação} a menos de isomorfismo
--- ficam determinadas pelo \omterm{def:b3:representations:character}{caráter}. Além disso, $\langle \chi_V, \chi_V\rangle = \sum_i
m_i^2$; em particular,
$V$ é \omterm{def:b3:representations:rep}{irredutível} se, e somente se, $\langle\chi_V,
\chi_V\rangle = 1$.
\end{corollary}

\begin{proof}
$\chi_V = \sum m_i\chi_{V_i}$ (\cref{prop:b3:representations:charbasics}); tome produtos internos com cada $\chi_{V_i}$ e use a
ortonormalidade. Duas \omterm{def:b3:representations:rep}{representações} de \omterm{def:b3:representations:character}{caracteres} iguais têm
multiplicidades iguais e, portanto, são isomorfas.
\end{proof}

\begin{theorem}[A representação regular]
\label{thm:b3:representations:regular}
Sejam $\chi_1, \dots, \chi_r$ os \omterm{def:b3:representations:character}{caracteres} irredutíveis distintos, de graus $n_i = \chi_i(e)$. A
\omterm{ex:b3:representations:examples}{representação regular} se decompõe com multiplicidades $m_i = n_i$;
consequentemente,
\[
\sum_{i=1}^{r} n_i^2 = \abs G,
\qquad
\sum_i n_i\chi_i(g) = 0 \quad (g \neq e).
\]
\end{theorem}

\begin{proof}
O \omterm{def:b3:representations:character}{caráter} regular: $\chi_{\mathrm{reg}}(g) =
\#\{h : gh = h\}$, que vale $\abs G$ para $g = e$ e
$0$ caso contrário. Logo $m_i = \langle \chi_i,\chi_{\mathrm{reg}}\rangle =
\frac1{\abs G}\overline{\chi_i(e)}\,\abs G = n_i$. Avaliando $\chi_{\mathrm{reg}} = \sum n_i\chi_i$ em $e$ e em
$g \neq e$ obtêm-se as duas identidades exibidas.
\end{proof}

\begin{theorem}\label{thm:b3:representations:numberirr}
Os \omterm{def:b3:representations:character}{caracteres} irredutíveis formam uma \emph{base} ortonormal de $\mathcal{CF}(G)$:
o número de \omterm{def:b3:representations:rep}{representações irredutíveis} é igual ao número de \omterm{ex:b3:groups:actions}{classes de
conjugação} de $G$.
\end{theorem}

\begin{proof}
Falta apenas a completude: seja $f \in \mathcal{CF}(G)$ ortogonal a todo $\chi_i$;
mostremos que $f = 0$. Para uma \omterm{def:b3:representations:rep}{representação} $(V,\rho)$, ponha $T_{f,\rho} = \frac{1}{\abs G}
\sum_g \overline{f(g)}\,\rho(g)$.
Ela é equivariante: para $h \in
G$,
\[
\rho(h)T_{f,\rho}\rho(h)^{-1} = \frac1{\abs G}\sum_g
\overline{f(g)}\rho(hgh^{-1})
= \frac1{\abs G}\sum_{g'}\overline{f(h^{-1}g'h)}\rho(g') =
T_{f,\rho}
\]
($f$ é \omterm{def:b3:representations:character}{função de classe}). Se $\rho$ é \omterm{def:b3:representations:rep}{irredutível} de grau
$n$, Schur dá $T_{f,\rho} = \lambda\,\mathrm{id}$ com
\[
\lambda = \frac{\operatorname{tr}T_{f,\rho}}{n}
= \frac{1}{n\abs G}\sum_g \overline{f(g)}\,\chi_\rho(g)
= \frac1n\,\langle f, \chi_\rho\rangle = 0 .
\]
Logo $T_{f,\rho} = 0$ em toda \omterm{def:b3:representations:rep}{irredutível} e, portanto (somas diretas), em
\emph{toda} \omterm{def:b3:representations:rep}{representação} --- em particular na regular. Aplique-a ao
vetor da base $e_e$: $0 = T_{f,\mathrm{reg}}e_e =
\frac1{\abs G}\sum_g\overline{f(g)}e_g$, o que força todo $\overline{f(g)} = 0$. Assim, a família
ortonormal $(\chi_i)$ gera $\mathcal{CF}(G)$, cuja dimensão é o número de \omterm{ex:b3:groups:actions}{classes de
conjugação}.
\end{proof}

\begin{corollary}[Ortogonalidade das colunas]
\label{cor:b3:representations:column}
Para $g, h \in G$:
\[
\sum_{i=1}^{r}\overline{\chi_i(g)}\,\chi_i(h) =
\begin{cases}
\abs{Z_G(g)} & \text{se } g, h \text{ são conjugados},\\
0 & \text{caso contrário.}
\end{cases}
\]
\end{corollary}

\begin{proof}
Sejam $g_1, \dots, g_r$ representantes das classes e $c_j$ os tamanhos das
classes. A matriz $U_{ij} =
\sqrt{c_j/\abs G}\;\chi_i(g_j)$ de tamanho $r \times r$ tem linhas ortonormais
(\cref{thm:b3:representations:orthogonality} escrito por classes: $\sum_j \frac{c_j}{\abs G}\chi_i(g_j)\overline{\chi_{i'}(g_j)} =
\delta_{ii'}$), isto é, $UU^* = I$; uma matriz
quadrada com $UU^* =
I$ satisfaz também $U^*U = I$: as colunas são ortonormais, o
que se desdobra na identidade exibida ($\abs G/c_j = \abs{Z_G(g_j)}$, a relação
órbita--estabilizador para a conjugação).
\end{proof}

\begin{proposition}[Caracteres unidimensionais; levantamento]
\label{prop:b3:representations:onedim}
(a) $G$ é abeliano se, e somente se, todas as suas \omterm{def:b3:representations:rep}{representações
irredutíveis} têm grau $1$; o número de \omterm{def:b3:representations:character}{caracteres} de grau $1$ de
um $G$ qualquer é $[G : D(G)]$ (eles são os \omterm{def:b3:representations:character}{caracteres} da abelianização).
(b) Se $N \trianglelefteq G$, as \omterm{def:b3:representations:rep}{representações irredutíveis} de $G/N$ se levantam
(compondo com $G \to G/N$) exatamente às \omterm{def:b3:representations:rep}{representações irredutíveis} de
$G$ cujo núcleo contém $N$.
\end{proposition}

\begin{proof}
(a) Se $G$ é abeliano, cada classe é unitária: $r = \abs G$, e $\sum n_i^2 = \abs G$ força
todos os $n_i = 1$; reciprocamente, se todos os $n_i = 1$, a \omterm{ex:b3:representations:examples}{representação
regular} é uma soma de \omterm{def:b3:representations:rep}{representações} unidimensionais, de modo que $\rho_{\mathrm{reg}}(G)$
é simultaneamente diagonalizável, logo comutativa, e $\rho_{\mathrm{reg}}$ é fiel:
$G$ é abeliano. As \omterm{def:b3:representations:rep}{representações} de grau $1$ são morfismos
$G \to \C^\times$ com alvo abeliano: elas se fatoram por $G^{\mathrm{ab}} = G/D(G)$
(\cref{exo:b3:groups:9}), e os \omterm{def:b3:representations:character}{caracteres} distintos do grupo abeliano $G^{\mathrm{ab}}$ são
em número de $\abs{G^{\mathrm{ab}}}$ (ele tem essa quantidade de classes, todas de grau
$1$). (b) A composição com a projeção preserva a irredutibilidade
(os subespaços invariantes se correspondem), e uma \omterm{def:b3:representations:rep}{representação}
trivial em $N$ se fatora pelo quociente (\cref{thm:b3:groups:firstiso}).
\end{proof}

\section{Tabelas de caracteres}

\begin{definition}\label{def:b3:representations:table}
A \emph{tabela de caracteres}\index{tabela de caracteres} de $G$ é
a matriz $\bigl(\chi_i(g_j)\bigr)$ de tamanho $r \times r$: linhas indexadas pelos \omterm{def:b3:representations:character}{caracteres}
irredutíveis, colunas pelas \omterm{ex:b3:groups:actions}{classes de conjugação} (com seus tamanhos
indicados). As linhas são ortonormais para o produto ponderado e as
colunas são ortogonais (\cref{cor:b3:representations:column}): a tabela é fortemente
sobredeterminada, e é isso que a torna calculável.
\end{definition}

\begin{example}[A tabela de $S_3$]\label{ex:b3:representations:s3}
Classes: $e$ (tamanho 1), transposições (3), $3$-ciclos (2);
logo $r
= 3$ irredutíveis, de graus $n_i$ com $\sum n_i^2 = 6$: $1, 1,
2$. Grau
$1$: a trivial $\mathbf 1$ e a assinatura $\varepsilon$. A última linha decorre da
ortogonalidade das colunas (ou de $\chi_{\mathrm{std}} = \chi_{\mathrm{perm}} - \mathbf 1$):
\[
\begin{array}{c|ccc}
S_3 & e & (1\,2)\ [3] & (1\,2\,3)\ [2]\\
\hline
\mathbf 1 & 1 & 1 & 1\\
\varepsilon & 1 & -1 & 1\\
\chi_{\mathrm{std}} & 2 & 0 & -1
\end{array}
\]
Verificação: $\langle\chi_{\mathrm{std}},\chi_{\mathrm{std}}\rangle =
\frac{1}{6}(4 + 0 + 2) = 1$: é \omterm{def:b3:representations:rep}{irredutível}.
\end{example}

\begin{example}[A tabela de $S_4$]\label{ex:b3:representations:s4}
Classes: $e$ [1], transposições [6], transposições duplas [3],
$3$-ciclos [8], $4$-ciclos [6]: cinco irredutíveis, $\sum n_i^2 =
24$ com
duas de grau $1$ ($\mathbf 1, \varepsilon$; $[S_4 :
D(S_4)] = [S_4 : A_4] = 2$): graus $1, 1, 2, 3, 3$. A de grau $2$ se
levanta de $S_4/V \cong S_3$ (\cref{prop:b3:representations:onedim}(b), com $V$ o grupo de Klein); as
de grau $3$: a \omterm{def:b3:representations:rep}{representação} padrão e sua \omterm{def:b3:modules:torsion}{torção} por $\varepsilon$:
\[
\begin{array}{c|ccccc}
S_4 & e\,[1] & (1\,2)\,[6] & (1\,2)(3\,4)\,[3] &
(1\,2\,3)\,[8] & (1\,2\,3\,4)\,[6]\\
\hline
\mathbf 1 & 1 & 1 & 1 & 1 & 1\\
\varepsilon & 1 & -1 & 1 & 1 & -1\\
\chi_2 & 2 & 0 & 2 & -1 & 0\\
\chi_{\mathrm{std}} & 3 & 1 & -1 & 0 & -1\\
\varepsilon\chi_{\mathrm{std}} & 3 & -1 & -1 & 0 & 1
\end{array}
\]
($\chi_{\mathrm{std}}(g) = \operatorname{fix}(g) - 1$; $\chi_2$ avalia a tabela de $S_3$ na imagem de cada
classe módulo $V$.) Todas as verificações de ortogonalidade de
linhas e colunas passam --- fazer duas delas é o aquecimento do~\cref{exo:b3:representations:3}.
\end{example}

\begin{method}\label{met:b3:representations:table}
Para construir uma \omterm{def:b3:representations:table}{tabela de caracteres}: (1) liste as \omterm{ex:b3:groups:actions}{classes de
conjugação} e seus tamanhos; (2) conte os \omterm{def:b3:representations:character}{caracteres} de grau $1$ por
meio de $G/D(G)$ e escreva-os; (3) encontre os graus restantes a
partir de $\sum n_i^2 = \abs G$ (combinatória de inteiros pequenos); (4) obtenha
irredutíveis baratas: levante de quocientes, subtraia $\mathbf 1$ dos
\omterm{def:b3:representations:character}{caracteres} de permutação (verifique $\langle\chi,\chi\rangle = 1$), multiplique \omterm{def:b3:representations:character}{caracteres} já
conhecidos por \omterm{def:b3:representations:character}{caracteres} de grau $1$; (5) termine as linhas
desconhecidas pela ortogonalidade das colunas --- cada coluna é
ortogonal às colunas já completas, e a coluna de $e$ carrega os
graus. Confira tudo com uma varredura completa de ortogonalidade.
\end{method}

\begin{omfigure}
\begin{tikzpicture}[scale=0.52]
  \draw[thick] (0,0) rectangle (6,6);
  \fill[omDef!25] (0,5) rectangle (1,6);
  \fill[omProp!25] (1,4) rectangle (2,5);
  \fill[omThm!20] (2,0) rectangle (6,4);
  \draw (0,5) rectangle (1,6);
  \draw (1,4) rectangle (2,5);
  \draw (2,0) rectangle (6,4);
  \node[font=\small] at (0.5,5.5) {$1$};
  \node[font=\small] at (1.5,4.5) {$1$};
  \node[font=\small] at (4,2) {$M_2(\C)$};
\end{tikzpicture}

{\small A \omterm{ex:b3:representations:examples}{representação regular} de $S_3$, diagonalizada por blocos:
$\C[S_3] \cong \C \times \C \times M_2(\C)$, de dimensões $1 + 1 +
4 = 6 = \abs{S_3}$. Em geral, $\C[G] \cong \prod_i
M_{n_i}(\C)$: a identidade $\sum n_i^2 = \abs G$ é uma
afirmação sobre blocos de matrizes.}
\end{omfigure}

\section{Exercícios}

\begin{exercise}[$\star$]\label{exo:b3:representations:1}
(a) Mostre que os \omterm{def:b3:representations:character}{caracteres} irredutíveis de $\Z/n\Z$ são os $\chi_k(\bar m) = \eu^{2\iu\pi km/n}$,
$k = 0, \dots, n-1$, e escreva a \omterm{def:b3:representations:table}{tabela de caracteres} de $\Z/4\Z$.
(b) Verifique nela as duas relações de ortogonalidade --- e reconheça a
matriz: onde este livro já a encontrou?
\end{exercise}

\begin{exercise}[$\star$]\label{exo:b3:representations:2}
Sejam $G$ agindo sobre um conjunto finito $X$ e $\chi$ o
\omterm{def:b3:representations:character}{caráter} da \omterm{def:b3:representations:rep}{representação} de permutação $\C^X$.
(a) Mostre que $\chi(g) = \abs{\operatorname{Fix}_X(g)}$ e que $\langle \mathbf 1, \chi\rangle = \#\{\text{órbitas}\}$ --- o lema de contagem de Burnside
(\cref{exo:b3:groups:5}) é um cálculo de \omterm{def:b3:representations:character}{caracteres}.
(b) Suponha a ação transitiva, de modo que $\chi = \mathbf 1 +
\psi$. Mostre que $\psi$
é \omterm{def:b3:representations:rep}{irredutível} se, e somente se, a ação é \emph{$2$-transitiva}
(transitiva nos pares ordenados de pontos distintos). \emph{(Calcule
$\langle\chi,\chi\rangle$ como o número de \omterm{def:b3:groups:action}{órbitas} em $X \times X$.)}
(c) Conclua que a \omterm{def:b3:representations:rep}{representação} padrão de $S_n$ ($n \geq
2$) é
\omterm{def:b3:representations:rep}{irredutível}.
\end{exercise}

\begin{exercise}[$\star$]\label{exo:b3:representations:3}
Reconstrua do zero a tabela de $S_3$ seguindo o~\cref{met:b3:representations:table} e
verifique em seguida duas relações de ortogonalidade de linhas e duas
de colunas na tabela de $S_4$ do~\cref{ex:b3:representations:s4}. Decomponha em
irredutíveis o \omterm{def:b3:representations:character}{caráter} de permutação de $S_4$ agindo sobre $\{1,2,3,4\}$ e
o \omterm{def:b3:representations:character}{caráter} $\chi_{\mathrm{std}}^2$ (quadrado ponto a ponto).
\end{exercise}

\begin{exercise}[$\star\star$]\label{exo:b3:representations:4}
Calcule as \omterm{def:b3:representations:table}{tabelas de caracteres} de $D_4$ e de $Q_8$. Conclua
que dois grupos não isomorfos podem ter \omterm{def:b3:representations:table}{tabelas de caracteres} idênticas
--- que dados de teoria de grupos a tabela ainda assim captura nesse
par (ordens dos \omterm{ex:b3:groups:actions}{centros}, abelianizações, número de involuções)? Quais
deles ela \emph{deixa} de capturar?
\end{exercise}

\begin{exercise}[$\star\star$]\label{exo:b3:representations:5}
\omterm{def:b3:representations:table}{Tabela de caracteres} de $A_4$: classes $e$ [1], transposições
duplas [3] e \emph{duas} classes de $3$-ciclos [4], [4].
(a) Explique a cisão dos $3$-ciclos (compare os \omterm{ex:b3:groups:actions}{centralizadores} em
$S_4$ e em $A_4$, como no~\cref{exo:b3:groups:11}).
(b) Encontre os três \omterm{def:b3:representations:character}{caracteres} de grau $1$ (via $A_4/V \cong
\Z/3\Z$) e o \omterm{def:b3:representations:character}{caráter}
de grau $3$ (restrinja $\chi_{\mathrm{std}}$ de $S_4$), e monte a tabela.
(c) Leia na tabela os \omterm{def:b3:groups:normal}{subgrupos normais} de $A_4$ (núcleos $\{g : \chi(g) = \chi(e)\}$ e
suas interseções).
\end{exercise}

\begin{exercise}[$\star\star$]\label{exo:b3:representations:6}
(a) Mostre que $\ker\chi = \{g : \chi(g) = \chi(e)\}$ é o núcleo da \omterm{def:b3:representations:rep}{representação} subjacente
(\cref{prop:b3:representations:charbasics}, caso de igualdade).
(b) Mostre que todo \omterm{def:b3:groups:normal}{subgrupo normal} de $G$ é uma interseção de
núcleos de \omterm{def:b3:representations:character}{caracteres} irredutíveis. \emph{(Represente $G/N$
fielmente: sua \omterm{ex:b3:representations:examples}{representação regular}.)}
(c) Deduza: $G$ é simples se, e somente se, $\ker\chi_i = \{e\}$ para todo
$\chi_i$ \omterm{def:b3:representations:rep}{irredutível} não trivial --- a simplicidade é legível na
\omterm{def:b3:representations:table}{tabela de caracteres}.
\end{exercise}

\begin{exercise}[$\star\star$]\label{exo:b3:representations:7}
Graus de $\abs G = 8$: mostre que um grupo não abeliano de ordem $8$ tem
padrão de graus $(1,1,1,1,2)$ e que sua \omterm{def:b3:representations:rep}{representação} de grau $2$ é fiel.
Mais geralmente, mostre que um grupo não abeliano de ordem $p^3$ tem
padrão $(1^{\,p^2},
p, \dots, p)$, com $p^2$ uns e $p - 1$ \omterm{def:b3:representations:character}{caracteres} de grau
$p$. \emph{(Use $[G : D(G)]$ e $\sum n_i^2 = \abs G$; aqui $D(G) = Z(G)$ tem ordem $p$.)}
\end{exercise}

\begin{exercise}[$\star\star\star$]\label{exo:b3:representations:8}
Para grupos finitos $G, H$: mostre que as \omterm{def:b3:representations:character}{funções de classe} $\chi(g)\psi(h)$
em $G \times H$, com $\chi, \psi$ \omterm{def:b3:representations:character}{caracteres} irredutíveis de $G, H$, são
exatamente os \omterm{def:b3:representations:character}{caracteres} irredutíveis de $G \times H$. \emph{(A ortonormalidade
é um cálculo direto; a completude, uma contagem de classes.)} Deduza a
\omterm{def:b3:representations:table}{tabela de caracteres} de $\Z/2\Z \times \Z/2\Z$ e redemonstre a~\cref{prop:b3:representations:onedim}(a) para
grupos abelianos finitos via o teorema de estrutura.
\end{exercise}

\begin{exercise}[$\star\star\star$]\label{exo:b3:representations:9}
A \omterm{def:b3:representations:table}{tabela de caracteres} de $A_5$ (classes de tamanhos $1, 15, 20, 12,
12$, do~\cref{exo:b3:groups:11}):
(a) Mostre que os graus são $1, 3, 3, 4, 5$ \emph{(a única solução de $\sum n_i^2 = 60$ com
$n_1 = 1$ e, usando o~\cref{exo:b3:representations:6}(c) com a simplicidade, nenhum outro
$n_i = 1$)}.
(b) Construa o \omterm{def:b3:representations:character}{caráter} de grau $4$ (ação por permutação em $5$
pontos) e o de grau $5$ (a ação sobre os seis $5$-subgrupos de
Sylow dá grau $6 = 1 + 5$; verifique a irredutibilidade), e complete as duas
linhas de grau $3$ pela ortogonalidade das colunas: entradas com a
razão áurea $\frac{1\pm\sqrt5}2$ aparecem nas classes dos $5$-ciclos.
(c) Verifique, a partir da tabela pronta, que $A_5$ é simples
(\cref{exo:b3:representations:6}(c)).
\end{exercise}

\begin{exercise}[$\star\star$]\label{exo:b3:representations:10}
Sejam $\rho$ uma \omterm{def:b3:representations:rep}{representação irredutível} de grau $n$ e $z
\in Z(G)$.
Mostre que $\rho(z) = \lambda_z\,\mathrm{id}$ com $\lambda \colon Z(G) \to \C^\times$ um morfismo (o \emph{\omterm{def:b3:representations:character}{caráter} central}) e
deduza $\abs{\chi(z)} = n$ para $z$ central. Aplicação: se $G$ tem uma
\omterm{def:b3:representations:rep}{representação irredutível} fiel, então $Z(G)$ é cíclico.
\end{exercise}

\begin{exercise}[$\star\star$]\label{exo:b3:representations:11}
(Projeções isotípicas) Sejam $(V, \rho)$ uma \omterm{def:b3:representations:rep}{representação} de $G$ e
$\chi_i$ um \omterm{def:b3:representations:character}{caráter} \omterm{def:b3:representations:rep}{irredutível} de grau $n_i$. Defina
\[
p_i = \frac{n_i}{\abs G}\sum_{g\in G}
\overline{\chi_i(g)}\,\rho(g) \;\in\; \mathcal L(V).
\]
(a) Mostre que $p_i$ é $G$-equivariante e calcule sua restrição
a uma subrepresentação \omterm{def:b3:representations:rep}{irredutível} $W \subseteq
V$ de \omterm{def:b3:representations:character}{caráter} $\chi_j$: ela é
$\delta_{ij}\,
\mathrm{id}_W$ \emph{(Schur; tome traços para identificar o escalar)}.
(b) Deduza que $p_i$ é uma projeção sobre a soma $V_i$ de todas
as subrepresentações irredutíveis de \omterm{def:b3:representations:character}{caráter} $\chi_i$ (a
\emph{componente isotípica}), que $\sum_ip_i =
\mathrm{id}_V$, e que a decomposição $V =
\bigoplus_iV_i$ é
canônica --- ao contrário da divisão mais fina de cada $V_i$ em
irredutíveis.
(c) Para a \omterm{ex:b3:representations:examples}{representação regular} de $S_3$ e o \omterm{def:b3:representations:character}{caráter} assinatura
$\varepsilon$, escreva $p_\varepsilon$ explicitamente como elemento da álgebra do grupo e
verifique $p_\varepsilon^2
= p_\varepsilon$ à mão.
\end{exercise}

\begin{exercise}[$\star\star$]\label{exo:b3:representations:12}
(Lendo uma tabela) A \omterm{def:b3:representations:table}{tabela de caracteres} de um certo grupo $G$ de
ordem $24$ é parcialmente conhecida: ele tem $5$ classes, de
tamanhos $1, 6, 8, 6, 3$, e graus $1, 1, 2, 3, 3$.
(a) Recupere a tabela completa: os dois \omterm{def:b3:representations:character}{caracteres} lineares (um
trivial; o outro assume o valor $-1$ exatamente nas classes de
tamanhos $6$ e $6$), depois o \omterm{def:b3:representations:character}{caráter} de grau $2$ pela
ortogonalidade com a coluna da identidade e, do mesmo modo, os dois
\omterm{def:b3:representations:character}{caracteres} de grau $3$ (um deles é $\chi_2\chi_4$).
(b) Identifique $G$ ($\cong S_4$: compare as classes com os tipos de
ciclo) e extraia da tabela os \omterm{def:b3:groups:normal}{subgrupos normais} via o~\cref{exo:b3:representations:6}:
os núcleos de $\chi_2$ (índice $2$: $A_4$) e de
$\chi_3$ (o grupo de Klein $V_4$), e nada mais além de $\{e\}, G$.
(c) Explique como a tabela mostra que $G/V_4 \cong S_3$ \emph{(quais \omterm{def:b3:representations:character}{caracteres} se
fatoram pelo quociente?)}.
\end{exercise}

\section{Problema: o teorema \texorpdfstring{$p^aq^b$}{paqb} de
Burnside}

\begin{problem}[{Problema de fim de semana --- solubilidade dos grupos
de ordem $p^aq^b$}]\label{pb:b3:representations:1}
Burnside demonstrou em 1904 que todo grupo cuja ordem tem no máximo
dois fatores primos é \omterm{def:b3:groups:derived}{solúvel} --- uma afirmação sobre grupos abstratos
cujas únicas demonstrações conhecidas durante meio século passavam pela
teoria de \omterm{def:b3:representations:character}{caracteres}. Este problema constrói a demonstração por
inteiro, reunindo o~\cref{ch:b3:groups} (\omterm{def:b3:groups:derived}{solubilidade}), o~\cref{ch:b3:modules}
($\Z$-módulos \omterm{def:b3:modules:free}{finitamente gerados}) e este capítulo. Ao longo de todo
o problema, $\chi_1, \dots, \chi_r$ são os \omterm{def:b3:representations:character}{caracteres} irredutíveis de $G$, $n_i = \chi_i(e)$.

\medskip
\noindent\textbf{Parte I --- Inteiros \omterm{def:b3:galois:algebraic}{algébricos}.} Um \emph{inteiro
\omterm{def:b3:galois:algebraic}{algébrico}} é uma raiz de um polinômio \emph{mônico} de $\Z[X]$.
\begin{enumerate}
  \item Mostre que $\alpha$ é um inteiro \omterm{def:b3:galois:algebraic}{algébrico} se, e somente
        se, o anel $\Z[\alpha]$ é um $\Z$-módulo \omterm{def:b3:modules:free}{finitamente gerado}.
  \item Deduza que os inteiros \omterm{def:b3:galois:algebraic}{algébricos} formam um subanel de
        $\C$. \emph{(Se $\Z[\alpha], \Z[\beta]$ são \omterm{def:b3:modules:free}{finitamente gerados}, $\Z[\alpha, \beta]$ também
        é, e os submódulos de $\Z$-módulos \omterm{def:b3:modules:free}{finitamente gerados} são
        \omterm{def:b3:modules:free}{finitamente gerados}, pelo~\cref{thm:b3:modules:submodule} e por um argumento de
        apresentação --- ou diretamente: um submódulo de $\Z^n$ é
        livre de posto $\leq n$.)}
  \item Mostre que um inteiro \omterm{def:b3:galois:algebraic}{algébrico} \emph{racional} é um inteiro.
        \emph{(Teorema da raiz racional.)}
  \item Mostre que todo valor de \omterm{def:b3:representations:character}{caráter} $\chi(g)$ é um inteiro \omterm{def:b3:galois:algebraic}{algébrico}.
\end{enumerate}

\medskip
\noindent\textbf{Parte II --- As relações das somas de classe.} Fixe
uma \omterm{def:b3:representations:rep}{representação irredutível} $(V, \rho)$ de grau $n$ e \omterm{def:b3:representations:character}{caráter}
$\chi$. Para uma \omterm{ex:b3:groups:actions}{classe de conjugação} $C$, ponha $S_C = \sum_{g \in C}\rho(g) \in
\mathcal L(V)$.
\begin{enumerate}[resume]
  \item Mostre que $S_C$ é equivariante, logo $S_C =
        \omega_C\,\mathrm{id}$ com
        \[
        \omega_C = \frac{\abs C\,\chi(g_C)}{n}
        \qquad (g_C \in C \text{ qualquer representante}).
        \]
  \item Mostre que $S_CS_{C'} = \sum_{C''}
        a_{CC'C''}\,S_{C''}$, em que $a_{CC'C''} \in \N$ conta, para um $z \in C''$ fixado, os
        pares $(x, y) \in C \times
        C'$ com $xy = z$. Deduza que os $\omega_C$ satisfazem
        $\omega_C\,\omega_{C'} = \sum_{C''}
        a_{CC'C''}\,\omega_{C''}$.
  \item Conclua que cada $\omega_C$ é um inteiro \omterm{def:b3:galois:algebraic}{algébrico}. \emph{(O
        $\Z$-módulo gerado por $1$ e pelos $\omega_C$ é um anel
        \omterm{def:b3:modules:free}{finitamente gerado}; aplique o critério da questão 1 --- mais
        precisamente, mostre que $M = \Z\text{-gerado}
        (1, (\omega_C)_C)$ satisfaz $\omega_{C_0} M \subseteq
        M$ e use um truque de
        determinante/Cayley--Hamilton, ou o argumento de subanel da
        questão 2.)}
  \item Deduza a \emph{divisibilidade de \omterm{thm:b3:galois:finitefields}{Frobenius}}: $n_i$ divide
        $\abs G$ para todo grau \omterm{def:b3:representations:rep}{irredutível}. \emph{(Calcule $\frac{\abs G}{n} = \frac{\abs G}{n}
        \langle\chi,\chi\rangle = \sum_C \omega_C\,
        \overline{\chi(g_C)}$:
        um inteiro \omterm{def:b3:galois:algebraic}{algébrico} que é racional.)}
\end{enumerate}

\medskip
\noindent\textbf{Parte III --- O critério de simplicidade de
Burnside.}
\begin{enumerate}[resume]
  \item Sejam $\chi$ \omterm{def:b3:representations:rep}{irredutível} de grau $n$ e $C$ uma classe
        com $\gcd(\abs C, n) = 1$. Usando Bézout e as questões 4--7, mostre que $\frac{\chi(g_C)}{n}$ é
        um inteiro \omterm{def:b3:galois:algebraic}{algébrico}.
  \item Suponha, além disso, $0 < \abs{\chi(g_C)} < n$. Mostre que isso é impossível: o
        inteiro \omterm{def:b3:galois:algebraic}{algébrico} $\alpha =
        \chi(g_C)/n$ tem todos os seus \emph{conjugados} ---
        os números $\alpha_\sigma = \frac1n\sigma(\chi(g_C))$ para $\sigma \in \operatorname{Gal}(\Q(\zeta_{\abs G})/\Q)$, cada um uma média de $n$
        raízes da unidade --- de módulo $\leq 1$, de modo que o
        produto $N = \prod_\sigma\alpha_\sigma$ é um inteiro \omterm{def:b3:galois:algebraic}{algébrico} racional com $0 < \abs N < 1$ ---
        justifique cada afirmação, citando o~\cref{thm:b3:galois:cyclotomicirred} para o \omterm{def:b3:galois:galois}{grupo
        de Galois} e o fato de que $\sigma$ permuta raízes da
        unidade. Conclua: ou $\chi(g_C) = 0$, ou $\rho(g_C)$ é escalar
        (\cref{prop:b3:representations:charbasics}).
  \item (Critério de Burnside) Seja $C \neq \{e\}$ uma \omterm{ex:b3:groups:actions}{classe de conjugação} de
        tamanho \emph{potência de primo} $p^k > 1$, e suponha $G$
        simples não abeliano. A ortogonalidade da coluna de $C$
        contra a coluna de $e$ dá $1 + \sum_{i \geq 2} n_i\chi_i(g_C) = 0$. Mostre que algum
        $\chi_i$ não trivial com $p \nmid n_i$ satisfaz $\chi_i(g_C) \neq 0$ \emph{(do
        contrário, $\frac1p$ seria um inteiro \omterm{def:b3:galois:algebraic}{algébrico})}; pela questão 10,
        $\rho_i(g_C)$ é escalar; obtenha uma contradição com a simplicidade
        \emph{(o conjunto dos $g$ com $\rho_i(g)$ escalar é um \omterm{def:b3:groups:normal}{subgrupo
        normal}; use a fidelidade, do~\cref{exo:b3:representations:6})}. Conclua:
        \emph{nenhum \omterm{def:b3:groups:simple}{grupo simples} não abeliano tem uma \omterm{ex:b3:groups:actions}{classe de
        conjugação} de tamanho potência de primo $> 1$.}
\end{enumerate}

\medskip
\noindent\textbf{Parte IV --- O teorema.}
\begin{enumerate}[resume]
  \item Seja $\abs G = p^aq^b$ com $a + b \geq 1$. Se $G$ é simples, mostre que ele é
        abeliano: escolha $z \neq e$ no \omterm{ex:b3:groups:actions}{centro} de um $q$-subgrupo de
        Sylow (\cref{thm:b3:groups:pfixed}) e considere o tamanho de sua \omterm{ex:b3:groups:actions}{classe de
        conjugação} $[G : Z_G(z)]$, uma potência de $p$ (por quê?); aplique a
        questão 11.
  \item Conclua por indução sobre $\abs G$: \emph{todo grupo de
        ordem $p^aq^b$ é \omterm{def:b3:groups:derived}{solúvel}} (Burnside). Por que o argumento
        se quebra para três primos --- e por que ele tem de se quebrar,
        dado que $\abs{A_5} = 2^2\cdot3\cdot5$?
\end{enumerate}

\medskip
\noindent\textbf{Parte V --- A \omterm{def:b3:representations:table}{tabela de caracteres} de $A_5$.} O
menor grupo que o teorema de Burnside não alcança merece seu retrato
completo; tudo o que vem abaixo usa apenas este capítulo e o~\cref{exo:b3:groups:11}.
\begin{enumerate}[resume]
  \item Recorde do~\cref{exo:b3:groups:11} as cinco \omterm{ex:b3:groups:actions}{classes de conjugação} de
        $A_5$: $\{e\}$, as $15$ transposições duplas, os
        $20$ ciclos de comprimento três e duas classes de $12$
        ciclos de comprimento cinco cada, representadas por $c =
        (1\,2\,3\,4\,5)$ e
        $c^2$. Explique por que os ciclos de comprimento cinco se
        cindem em duas $A_5$-classes, embora formem uma única
        $S_5$-classe.
  \item Mostre que os graus irredutíveis de $A_5$ são exatamente
        $1, 3, 3, 4, 5$: use $\sum_in_i^2 = 60$ com $r = 5$ classes e o fato de que
        $A_5$ é \omterm{prop:b3:galois:perfect}{perfeito} ($D(A_5) =
        A_5$), de modo que o \omterm{def:b3:representations:character}{caráter} trivial é
        seu único \omterm{def:b3:representations:character}{caráter} linear; depois elimine todo outro
        multiconjunto \emph{(escreva $59$ como soma de quatro
        quadrados de inteiros $\geq 2$: verifique que isso ocorre de
        um único modo com todas as parcelas sendo graus plausíveis)}.
  \item Seja $\pi$ o \omterm{def:b3:representations:character}{caráter} de permutação de $A_5$ em $\{1, \dots, 5\}$:
        $\pi(g) = \#\operatorname{Fix}(g)$, de valores $5, 1, 2, 0, 0$ nas cinco classes. Calcule $\langle\pi, \mathbf 1\rangle$ e
        $\langle\pi,
        \pi\rangle$, e deduza que $\chi_4 = \pi - \mathbf 1$ é \omterm{def:b3:representations:rep}{irredutível} de grau $4$, de
        valores $4, 0, 1, -1,
        -1$.
  \item Mesmo jogo nos $10$ pares não ordenados $\{i, j\}$: as contagens
        de pontos fixos são $10, 2, 1, 0, 0$. Calcule $\langle\pi_{10}, \pi_{10}\rangle$, $\langle\pi_{10},
        \mathbf 1\rangle$ e $\langle\pi_{10}, \chi_4\rangle$, deduza
        a decomposição $\pi_{10} = \mathbf 1 + \chi_4
        + \chi_5$ e obtenha o \omterm{def:b3:representations:character}{caráter} \omterm{def:b3:representations:rep}{irredutível}
        $\chi_5$ de grau $5$, de valores $5, 1, -1, 0, 0$.
  \item As duas irredutíveis restantes $\chi_2, \chi_3$ têm grau $3$. A
        ortogonalidade das colunas (cada coluna não identidade contra a
        coluna da identidade, e cada coluna consigo mesma) determina
        seus valores fora dos ciclos de comprimento cinco: mostre que
        $\chi_2(g) = \chi_3(g) = -1$ nas transposições duplas e $0$ nos ciclos de
        comprimento três.
  \item Nas classes dos ciclos de comprimento cinco, ponha $x = \chi_2(c)$ e
        $y =
        \chi_2(c^2)$; a simetria permite tomar $\chi_3(c) = y$, $\chi_3(c^2) = x$. Da coluna de
        $c$ emparelhada com a coluna da identidade e com a coluna de
        $c^2$, deduza $x + y = 1$ e $xy = -1$ (e confira o valor $x^2 +
        y^2 = 3$ dado
        pela coluna de $c$ consigo mesma), donde
        \[
        \{x, y\} = \Bigl\{\frac{1 + \sqrt5}2,\ \frac{1 -
        \sqrt5}2\Bigr\} :
        \]
        a razão áurea e sua conjugada. Monte a \omterm{def:b3:representations:table}{tabela de caracteres}
        completa de $A_5$.
  \item Faça as verificações: a norma da linha de $\chi_2$ é
        $1$ (use $\varphi^2 + \bar\varphi^2 = 3$), $\langle\chi_2,
        \chi_3\rangle = 0$, e a divisibilidade de \omterm{thm:b3:galois:finitefields}{Frobenius}
        (questão 8) para os cinco graus. Onde, na tabela, você
        \emph{vê} uma diferença em relação a $S_5$, cujos valores de
        \omterm{def:b3:representations:character}{caráter} são todos inteiros racionais?
  \item Deduza apenas da tabela que $A_5$ é simples: um \omterm{def:b3:groups:normal}{subgrupo
        normal} é uma reunião de \omterm{ex:b3:groups:actions}{classes de conjugação} que contém
        $e$ e cuja cardinalidade divide $60$ --- verifique que
        nenhuma subsoma própria de $1 + 15 + 20 + 12 + 12$ contendo a parcela $1$
        divide $60$. Confronte com o critério da questão 11:
        verifique que nenhuma classe de $A_5$ tem tamanho potência
        de primo $> 1$.
  \item (Coda icosaédrica) $A_5$ é o grupo de rotações do
        icosaedro, e as \omterm{def:b3:representations:rep}{representações} de grau $3$ são as duas ações
        geométricas em $\R^3$. Verifique a identidade dos traços:
        uma rotação de ângulo $\theta$ tem traço $1 +
        2\cos\theta$, e $1 + 2\cos\frac{2\pi}5 =
        \frac{1+\sqrt5}2 = \varphi$.
        Explique sem nenhum cálculo por que o outro \omterm{def:b3:representations:character}{caráter} de grau
        $3$ tem de carregar o valor conjugado: o \omterm{def:b3:galois:galois}{grupo de Galois} de
        $\Q(\sqrt5)/\Q$ age sobre a \omterm{def:b3:representations:table}{tabela de caracteres} inteira (entrada a
        entrada), permutando os \omterm{def:b3:representations:character}{caracteres} irredutíveis.
\end{enumerate}

\medskip
\noindent\textbf{Parte VI --- Complementos: uma cota central e o
quadrado tensorial.}
\begin{enumerate}[resume]
  \item (Mais fino do que uma divisibilidade) Seja $\chi$
        \omterm{def:b3:representations:rep}{irredutível} de grau $n$. Mostre que $\abs{\chi(z)} = n$ para todo $z \in Z(G)$
        \emph{(lema de Schur: $\rho(z)$ é escalar, de ordem finita)} e
        deduza de $\langle\chi, \chi\rangle = 1$ a cota
        \[
        n^2 \leq [G : Z(G)] .
        \]
        Mostre que os grupos não abelianos de ordem $8$ têm graus
        irredutíveis $1, 1, 1, 1, 2$ \emph{(cinco \omterm{ex:b3:groups:actions}{classes de conjugação}; escreva
        $8$ como soma de cinco quadrados)} e atingem a igualdade
        $4 = [G : Z(G)]$; verifique a cota em $A_5$, cujo \omterm{ex:b3:groups:actions}{centro} é trivial.
  \item (Quadrado tensorial de $\chi_2$) Para $g$ de ordem
        finita, $\rho(g)$ é diagonalizável com autovalores raízes da
        unidade; deduza as fórmulas de \omterm{def:b3:representations:character}{caráter}
        \[
        \chi_{\operatorname{Sym}^2 V}(g)
        = \frac{\chi(g)^2 + \chi(g^2)}2,
        \qquad
        \chi_{\Lambda^2 V}(g)
        = \frac{\chi(g)^2 - \chi(g^2)}2 .
        \]
        Aplique-as a $\chi_2$ de $A_5$ \emph{(note que $g^2$
        percorre a classe de $c^2$ quando $g$ percorre a de
        $c$, e reciprocamente)}: mostre que $\Lambda^2\chi_2 =
        \chi_2$ e $\operatorname{Sym}^2\chi_2 = \mathbf 1 +
        \chi_5$, donde
        \[
        \chi_2\otimes\chi_2 = \mathbf 1 + \chi_2 + \chi_5 .
        \]
        Interprete $\Lambda^2\chi_2 = \chi_2$ geometricamente por meio do produto vetorial
        em $\R^3$.
  \item (Auditoria final da tabela) Verifique numericamente: a
        ortogonalidade entre as duas colunas dos ciclos de comprimento
        cinco ($\varphi\bar\varphi = -1$), o valor $5 =
        \abs{Z_{A_5}(c)}$ para a coluna de $c$ consigo
        mesma, e o anulamento do \omterm{def:b3:representations:character}{caráter} regular $\sum_i
        n_i\chi_i(g) = 0$ em cada uma das
        quatro colunas não identidade da tabela.
\end{enumerate}
\end{problem}
